
\documentclass[tikz]{standalone}
\usepackage{tikz}
\usepackage{pgfplots}
\usepackage{amsmath}
\pgfplotsset{compat=1.17}

\begin{document}

% Frame 1: Draw axes
\begin{tikzpicture}
  \begin{axis}[
    axis lines=middle,
    xmin=-5, xmax=5,
    ymin=-5, ymax=5,
    xlabel={$x$}, ylabel={$y$},
    ticks=none,
    width=6cm, height=6cm,
    domain=-5:5,
  ]
  \end{axis}
\end{tikzpicture}

% Frame 2: Add first point
\begin{tikzpicture}
  \begin{axis}[
    axis lines=middle,
    xmin=-5, xmax=5,
    ymin=-5, ymax=5,
    xlabel={$x$}, ylabel={$y$},
    ticks=none,
    width=6cm, height=6cm,
    domain=-5:5,
  ]
    \addplot[only marks, red] coordinates {(1,0)};
    \node[above right] at (axis cs:-1,0) {\tiny$R_1(1;0)$};
  \end{axis}
\end{tikzpicture}

% Frame 3: Add second point
\begin{tikzpicture}
  \begin{axis}[
    axis lines=middle,
    xmin=-5, xmax=5,
    ymin=-5, ymax=5,
    xlabel={$x$}, ylabel={$y$},
    ticks=none,
    width=6cm, height=6cm,
    domain=-5:5,
  ]
    \addplot[only marks, red] coordinates {(1,0) (2,0)};
    \node[above right] at (axis cs:-1,0) {\tiny $R_1(1;0)$};
    \node[above right] at (axis cs:2,0) {\tiny$R_2(2;0)$};
  \end{axis}
\end{tikzpicture}

% Frame 4: Add third point
\begin{tikzpicture}
  \begin{axis}[
    axis lines=middle,
    xmin=-5, xmax=5,
    ymin=-5, ymax=5,
    xlabel={$x$}, ylabel={$y$},
    ticks=none,
    width=6cm, height=6cm,
    domain=-5:5,
  ]
    \addplot[only marks, red] coordinates {(1,0) (2,0) (0,-2)};
    \node[above right] at (axis cs:-1,0) {\tiny $R_1(1;0)$};
    \node[above right] at (axis cs:2,0) {\tiny $R_2(2;0)$};
    \node[above left] at (axis cs:0,-2) {\tiny $C(0;-2)$};
  \end{axis}
\end{tikzpicture}

% Frame 4: Add third point
\begin{tikzpicture}
  \begin{axis}[
    axis lines=middle,
    xmin=-5, xmax=5,
    ymin=-5, ymax=5,
    xlabel={$x$}, ylabel={$y$},
    ticks=none,
    width=6cm, height=6cm,
    domain=-5:5,
  ]
  \addplot[only marks, red] coordinates {(1,0) (2,0) (0,-2) (3/2,1/4)};
    \node[above right] at (axis cs:-1,0) {\tiny$R_1(1;0)$};
    \node[above right] at (axis cs:2,0) {\tiny$R_2(2;0)$};
    \node[above left] at (axis cs:0,-2) {\tiny$C(0;-2)$};
    \node[above] at (axis cs:3/2,1) {\tiny$S\left(\dfrac{3}{2};\dfrac{1}{4}\right)$};
  \end{axis}
\end{tikzpicture}

% Frames 5 to 14: Draw the parabola incrementally
\foreach \i in {-20,...,80} {
\begin{tikzpicture}
  \begin{axis}[
    axis lines=middle,
    xmin=-5, xmax=5,
    ymin=-5, ymax=5,
    xlabel={$x$}, ylabel={$y$},
    ticks=none,
    width=6cm, height=6cm,
    domain=-5:5,
  ]
  \addplot[only marks, red] coordinates {(1,0) (2,0) (0,-2) (3/2,1/4)};
    \node[above right] at (axis cs:-1,0) {\tiny $R_1(1;0)$};
    \node[above right] at (axis cs:2,0) {\tiny$R_2(2;0)$};
    \node[above left] at (axis cs:0,-2) {\tiny$C(0;-2)$};
    \node[above] at (axis cs:3/2,1) {\tiny$S\left(\dfrac{3}{2};\dfrac{1}{4}\right)$};

    \addplot[blue, thick, samples=100, domain=-5:5*\i/100] {-x^2 + 3*x - 2};
  \end{axis}
\end{tikzpicture}
}

\end{document}

